\documentclass[11pt,letterpaper]{article}

% --- Packages ---
\usepackage[margin=1in]{geometry}
\usepackage{graphicx}
\usepackage{booktabs}
\usepackage{longtable}
\usepackage{tabularx}
\usepackage{enumitem}
\usepackage{hyperref}
\usepackage{xcolor}
\usepackage{fancyhdr}
\usepackage{listings}
\usepackage{titlesec}
\usepackage[T1]{fontenc}
\usepackage{lmodern}
\usepackage{parskip}
\usepackage{array}
\usepackage{microtype}
\usepackage{tcolorbox}

% --- Colours ---
\definecolor{codebg}{gray}{0.95}
\definecolor{codeframe}{gray}{0.75}
\definecolor{linkblue}{HTML}{0066CC}
\definecolor{cautionred}{HTML}{CC0000}
\definecolor{notegray}{HTML}{555555}
\definecolor{guicolor}{HTML}{2E7D32}
\definecolor{tipgreen}{HTML}{1B5E20}
\definecolor{tipbg}{HTML}{E8F5E9}
\definecolor{warnbg}{HTML}{FFF3E0}
\definecolor{warnorange}{HTML}{E65100}
\definecolor{infobg}{HTML}{E3F2FD}
\definecolor{infoblue}{HTML}{0D47A1}

% --- Hyperlink styling ---
\hypersetup{
    colorlinks=true,
    linkcolor=linkblue,
    urlcolor=linkblue,
    citecolor=linkblue,
    pdftitle={X-CAVATE GUI User Guide},
    pdfauthor={X-CAVATE Team},
    pdfsubject={User Guide for the X-CAVATE Streamlit GUI},
}

% --- Code listing style ---
\lstset{
    backgroundcolor=\color{codebg},
    frame=single,
    rulecolor=\color{codeframe},
    basicstyle=\ttfamily\small,
    breaklines=true,
    breakatwhitespace=false,
    columns=fullflexible,
    keepspaces=true,
    showstringspaces=false,
    xleftmargin=4pt,
    xrightmargin=4pt,
    aboveskip=8pt,
    belowskip=8pt,
}

% --- Header / Footer ---
\pagestyle{fancy}
\fancyhf{}
\fancyhead[L]{\textbf{X-CAVATE GUI User Guide}}
\fancyhead[R]{\thepage}
\renewcommand{\headrulewidth}{0.4pt}

% --- Custom commands ---
\newcommand{\caution}[1]{\par\textcolor{cautionred}{\textbf{CAUTION:} #1}\par}
\newcommand{\note}[1]{\par\textcolor{notegray}{\textbf{Note:} #1}\par}
\newcommand{\code}[1]{\texttt{#1}}
\newcommand{\gui}[1]{\textcolor{guicolor}{\textsf{\textbf{#1}}}}
\newcommand{\btn}[1]{\fbox{\textsf{#1}}}
\newcolumntype{L}[1]{>{\raggedright\arraybackslash}p{#1}}
\newcolumntype{C}[1]{>{\centering\arraybackslash}p{#1}}
\tolerance=400
\emergencystretch=2em

% --- Callout boxes ---
\newtcolorbox{tipbox}{
    colback=tipbg,
    colframe=tipgreen,
    coltitle=white,
    fonttitle=\bfseries,
    title=Tip,
    boxrule=0.5pt,
    arc=2pt,
    left=6pt, right=6pt, top=4pt, bottom=4pt,
}

\newtcolorbox{warnbox}{
    colback=warnbg,
    colframe=warnorange,
    coltitle=white,
    fonttitle=\bfseries,
    title=Warning,
    boxrule=0.5pt,
    arc=2pt,
    left=6pt, right=6pt, top=4pt, bottom=4pt,
}

\newtcolorbox{infobox}{
    colback=infobg,
    colframe=infoblue,
    coltitle=white,
    fonttitle=\bfseries,
    title=Info,
    boxrule=0.5pt,
    arc=2pt,
    left=6pt, right=6pt, top=4pt, bottom=4pt,
}

% --- Title ---
\title{
    \vspace{-1cm}
    \textbf{\Huge X-CAVATE}\\[8pt]
    \LARGE GUI User Guide\\[6pt]
    \large Version 2.0.0
}
\author{}
\date{\today}

\begin{document}

% --- Title page ---
\begin{titlepage}
\centering
\vspace*{3cm}

{\Huge \textbf{X-CAVATE}}\\[12pt]
{\LARGE GUI User Guide}\\[20pt]
{\large Version 2.0.0}\\[40pt]

{\large Convert vascular network geometries into collision-free\\[4pt]
3D printer toolhead pathways and G-code.}

\vspace{2cm}

{\normalsize
\textbf{SM} = Single-Material \quad|\quad \textbf{MM} = Multimaterial
}

\vfill

{\normalsize \today}

\end{titlepage}

% --- Table of Contents ---
\tableofcontents
\newpage

% ===================================================================
\section{Getting Started}
% ===================================================================

X-CAVATE provides a web-based GUI built with Streamlit. It requires no coding
experience to operate. This section covers installation and launching the
interface.

\subsection{Installation}

\subsubsection{Option A: Conda (Recommended)}

\begin{enumerate}[leftmargin=*]
    \item Open a terminal (Mac: Terminal.app; Windows: Anaconda Prompt or PowerShell).
    \item Confirm conda is installed:
    \begin{lstlisting}
conda --version
    \end{lstlisting}
    If not installed, download Miniconda from \url{https://docs.conda.io/en/latest/miniconda.html}.
    \item Clone the repository and enter the directory:
    \begin{lstlisting}
git clone https://github.com/sohams-MASS/X-CAVATE.git
cd X-CAVATE
    \end{lstlisting}
    \item Create and activate the conda environment:
    \begin{lstlisting}
conda env create -f environment.yml
conda activate xcavate-gui
    \end{lstlisting}
    \item Verify the installation:
    \begin{lstlisting}
xcavate --help
    \end{lstlisting}
\end{enumerate}

\subsubsection{Option B: pip}

\begin{enumerate}[leftmargin=*]
    \item Confirm Python 3.9+ is installed: \code{python3 -{}-version}
    \item Clone the repository (same as above).
    \item Install with GUI support:
    \begin{lstlisting}
pip install ".[gui]"
    \end{lstlisting}
    For CLI-only usage (no web interface): \code{pip install .}
\end{enumerate}

\subsubsection{Option C: Docker}

\begin{enumerate}[leftmargin=*]
    \item Install Docker Desktop from \url{https://www.docker.com/products/docker-desktop}.
    \item Build the image:
    \begin{lstlisting}
docker build -t xcavate .
    \end{lstlisting}
    \item Run the container:
    \begin{lstlisting}
docker run -p 8501:8501 xcavate
    \end{lstlisting}
    \item Open \url{http://localhost:8501} in your browser.
\end{enumerate}

\subsection{Launching the GUI}

After installation (conda or pip), start the server:

\begin{lstlisting}
conda activate xcavate-gui      # if using conda
streamlit run xcavate/gui/app.py
\end{lstlisting}

A browser window opens automatically at \textbf{http://localhost:8501}. If it
does not open, navigate to that URL manually.

\begin{tipbox}
To stop the server, press \code{Ctrl+C} in the terminal.
\end{tipbox}

\subsection{Dependencies}

\begin{tabularx}{\textwidth}{l l X}
\toprule
\textbf{Package} & \textbf{Version} & \textbf{Purpose} \\
\midrule
numpy       & $\geq$ 1.21 & Array operations \\
pandas      & $\geq$ 1.3  & Data handling \\
scipy       & $\geq$ 1.7  & KD-tree spatial indexing, interpolation \\
plotly      & $\geq$ 5.0  & Interactive 3D visualization \\
matplotlib  & $\geq$ 3.5  & Static plots \\
streamlit   & $\geq$ 1.20 & Web GUI (optional with pip; included in conda env) \\
\bottomrule
\end{tabularx}

% ===================================================================
\section{GUI Overview}
\label{sec:overview}
% ===================================================================

The X-CAVATE GUI uses a two-panel layout:

\begin{description}[style=nextline]
    \item[\gui{Sidebar} (left panel)]
    File upload widgets and all configurable parameters in expandable sections.
    \item[\gui{Main area} (right panel)]
    Top-level tabs containing the pipeline controls, results, custom G-code
    editor, print instructions, and calibration validation.
\end{description}

\subsection{Abbreviations}

Throughout the interface:
\begin{itemize}[nosep]
    \item \textbf{SM} = Single-Material
    \item \textbf{MM} = Multimaterial
\end{itemize}

\subsection{Tab Structure}

The main area contains up to four tabs (the second appears only when
the Custom G-code toggle is enabled):

\begin{tabularx}{\textwidth}{l X}
\toprule
\textbf{Tab} & \textbf{Description} \\
\midrule
\gui{Pipeline}               & Run button, 7-step progress bar, 3D visualizations, download buttons \\
\gui{Custom G-code}          & Tabbed template editor for header, extrusion, pressure, and dwell G-code snippets (visible only when enabled) \\
\gui{Print Instructions}     & Step-by-step positioning and printing guidance for SM and MM modes \\
\gui{Calibration}            & Compute volumetric flow rate ($Q$) from calibration filaments and validate measured diameters against targets \\
\bottomrule
\end{tabularx}

% ===================================================================
\section{Sidebar: Input Files}
\label{sec:input-files}
% ===================================================================

The top of the sidebar contains two file uploaders:

\begin{tabularx}{\textwidth}{l l X}
\toprule
\textbf{Widget} & \textbf{Accepted type} & \textbf{Description} \\
\midrule
\gui{Network file} & \code{.txt} & Vascular network coordinates exported from SimVascular. Each vessel has a header line followed by coordinate rows (x, y, z in cm, plus optional radius and arterial/venous flag columns). \\
\gui{Inlet / Outlet file} & \code{.txt} & Text file identifying inlet and outlet node coordinates. Each section starts with the keyword \code{inlet} or \code{outlet}, followed by (x, y, z) triplets. \\
\bottomrule
\end{tabularx}

\begin{warnbox}
Both files must be uploaded before the \btn{Run X-CAVATE} button becomes active.
Coordinates in the network file must be in \textbf{centimeters} --- X-CAVATE
converts to millimeters internally.
\end{warnbox}

\subsection{Network File Format}

\begin{lstlisting}
Vessel: 0, Number of Points: 100
33.4878, 0.0000, 53.5303, 0.0200
33.4115, -0.5285, 52.7311, 0.0200
...
Vessel: 1, Number of Points: 100
15.0600, -37.3300, 11.6600, 0.0190
...
\end{lstlisting}

\begin{tabularx}{\textwidth}{l l X}
\toprule
\textbf{Column} & \textbf{Required} & \textbf{Description} \\
\midrule
1--3 & Yes & x, y, z coordinates (cm) \\
4    & No  & Vessel radius (needed for speed calculation) \\
5    & No  & Arterial/venous flag: 0 = venous, nonzero = arterial (needed for multimaterial) \\
\bottomrule
\end{tabularx}

\subsection{Inlet / Outlet File Format}

\begin{lstlisting}
inlet
21.97, -15.43, 50.00
11.81, 0.00, 21.34
outlet
36.05, -27.19, 50.00
50.00, 0.00, 26.54
\end{lstlisting}

Coordinates are matched to the nearest network point automatically.

% ===================================================================
\section{Sidebar: Required Parameters}
\label{sec:required-params}
% ===================================================================

These parameters appear directly below the file uploaders and must be set for
every run.

\begin{longtable}{L{3cm} l l L{6.5cm}}
\toprule
\textbf{Parameter} & \textbf{Type} & \textbf{Default} & \textbf{Description} \\
\midrule
\endfirsthead
\toprule
\textbf{Parameter} & \textbf{Type} & \textbf{Default} & \textbf{Description} \\
\midrule
\endhead
\gui{Nozzle diameter} (mm) & float & 0.5 & Outer diameter of the print nozzle. \\
\gui{Container height} (mm) & float & 10.0 & Height of the print container. \\
\gui{Decimal places} & int & 3 & Number of decimal places for output coordinate rounding. \\
\gui{Amount up} (mm) & float & 10.0 & Distance to raise the nozzle above the container between passes. \\
\gui{Printer type} & select & Pressure & Hardware target: Pressure (pneumatic extrusion), Positive Ink (displacement pump), or Aerotech (6-axis motion controller). \\
\gui{Multimaterial} & toggle & Off & Enable two-nozzle arterial/venous printing. Requires a 5-column input file (x, y, z, radius, arterial/venous flag). \\
\bottomrule
\end{longtable}

% ===================================================================
\section{Sidebar: Optional Parameters}
\label{sec:optional-params}
% ===================================================================

Optional parameters are organized into seven expandable sections. Click a
section header to reveal or hide its contents.

% --- 5.1 Tolerancing ---
\subsection{Tolerancing}

\begin{tabularx}{\textwidth}{L{3.2cm} l l X}
\toprule
\textbf{Parameter} & \textbf{Type} & \textbf{Default} & \textbf{Description} \\
\midrule
\gui{Enable tolerance} & toggle & Off & When on, nodes within the tolerance distance are allowed in the same pass even if one geometrically blocks the other. \\
\gui{Tolerance} & float & 0.0 & Tolerance distance (mm). 0 means none. \\
\bottomrule
\end{tabularx}

% --- 5.2 Print Settings ---
\subsection{Print Settings}

\begin{tabularx}{\textwidth}{L{3.2cm} l l X}
\toprule
\textbf{Parameter} & \textbf{Type} & \textbf{Default} & \textbf{Description} \\
\midrule
\gui{Print speed} (mm/s) & float & 1.0 & Nozzle travel speed during extrusion. Typical range: 0.5--5.0 mm/s. \\
\gui{Flow} (mm\textsuperscript{3}/s) & float & 0.161 & Volumetric flow rate for speed calculation mode. \\
\gui{Jog speed} (mm/s) & float & 5.0 & Speed for rapid non-printing travel moves between passes. \\
\gui{Jog speed lift} (mm/s) & float & 0.25 & Speed for the vertical lift at the end of each pass. \\
\gui{Initial lift} (mm) & float & 0.5 & Height the nozzle lifts after finishing a pass before traveling. \\
\gui{Dwell start} (s) & float & 0.08 & Pause after starting extrusion before nozzle movement. \textit{Aerotech only.} \\
\gui{Dwell end} (s) & float & 0.08 & Pause after nozzle stops before extrusion off. \textit{Aerotech only.} \\
\bottomrule
\end{tabularx}

\begin{infobox}
The \gui{Dwell start} and \gui{Dwell end} fields only appear when
the printer type is set to \textbf{Aerotech}.
\end{infobox}

% --- 5.3 Geometry ---
\subsection{Geometry}

\begin{tabularx}{\textwidth}{L{3.2cm} l l X}
\toprule
\textbf{Parameter} & \textbf{Type} & \textbf{Default} & \textbf{Description} \\
\midrule
\gui{Container X} (mm) & float & 50.0 & Width of the print container. Used for centering instructions. \\
\gui{Container Y} (mm) & float & 50.0 & Depth of the print container. Used for centering instructions. \\
\gui{Scale factor} & float & 1.0 & Multiplier applied to all coordinates after unit conversion. \\
\gui{Top padding} (mm) & float & 0.0 & Extra clearance above the highest network point for safe nozzle retraction. \\
\bottomrule
\end{tabularx}

% --- 5.4 Downsampling ---
\subsection{Downsampling}

\begin{tabularx}{\textwidth}{L{3.2cm} l l X}
\toprule
\textbf{Parameter} & \textbf{Type} & \textbf{Default} & \textbf{Description} \\
\midrule
\gui{Enable downsampling} & toggle & Off & Reduce point density to speed up printing. Branchpoints and endpoints are preserved. \\
\gui{Downsample factor} & int & 1 & Factor by which to downsample the network. \\
\bottomrule
\end{tabularx}

% --- 5.5 Multimaterial ---
\subsection{Multimaterial}

These settings are relevant when the \gui{Multimaterial} toggle is enabled.

\begin{tabularx}{\textwidth}{L{3.2cm} l l X}
\toprule
\textbf{Parameter} & \textbf{Type} & \textbf{Default} & \textbf{Description} \\
\midrule
\gui{Offset X} (mm) & float & 103.0 & Horizontal distance between the two printheads. \\
\gui{Offset Y} (mm) & float & 0.5 & Vertical fine-tuning offset between printheads. \\
\gui{Front nozzle} & int & 1 & 1 = venous nozzle in front, 2 = venous nozzle behind. \\
\gui{Printhead 1 name} & text & Aa & Axis designation for the arterial printhead in G-code. \\
\gui{Printhead 2 name} & text & Ab & Axis designation for the venous printhead in G-code. \\
\gui{Axis 1 name} & text & A & Z-axis letter for the arterial printhead. \\
\gui{Axis 2 name} & text & B & Z-axis letter for the venous printhead. \\
\gui{Jog translation} (mm/s) & float & 10.0 & Speed for travel moves between nozzles during printhead switching. \\
\gui{Resting pressure} (psi) & float & 0.0 & Pressure applied to the inactive printhead. \\
\gui{Active pressure} (psi) & float & 5.0 & Pressure applied to the active printhead during extrusion. \\
\bottomrule
\end{tabularx}

% --- 5.6 Positive Ink Displacement ---
\subsection{Positive Ink Displacement}

\begin{tabularx}{\textwidth}{L{3.5cm} l l X}
\toprule
\textbf{Parameter} & \textbf{Type} & \textbf{Default} & \textbf{Description} \\
\midrule
\gui{Use radii} & toggle & Off & Use per-node vessel radii instead of a fixed diameter. \\
\gui{Diameter} (mm) & float & 1.0 & Fixed vessel diameter for extrusion calculations. \\
\gui{Syringe diameter} (mm) & float & 1.0 & Inner diameter of the syringe barrel. \\
\gui{Factor} & float & 1.0 & Scaling multiplier for computed extrusion amounts. \\
\gui{Separate arterial/venous} & toggle & Off & Set different start/end displacements per material. \\
\bottomrule
\end{tabularx}

When \gui{Separate arterial/venous} is \textbf{off}, two fields appear:

\begin{tabularx}{\textwidth}{L{3.5cm} l l X}
\toprule
\textbf{Parameter} & \textbf{Type} & \textbf{Default} & \textbf{Description} \\
\midrule
\gui{Start} (mm) & float & 0.0 & Plunger displacement at the start of each pass to prime ink flow. \\
\gui{End} (mm) & float & 0.0 & Plunger displacement at the end of each pass (negative retracts). \\
\bottomrule
\end{tabularx}

When \gui{Separate arterial/venous} is \textbf{on}, four fields appear instead:

\begin{tabularx}{\textwidth}{L{3.5cm} l l X}
\toprule
\textbf{Parameter} & \textbf{Type} & \textbf{Default} & \textbf{Description} \\
\midrule
\gui{Start (arterial)} & float & 0.0 & Plunger displacement for the arterial syringe at pass start. \\
\gui{Start (venous)} & float & 0.0 & Plunger displacement for the venous syringe at pass start. \\
\gui{End (arterial)} & float & 0.0 & Plunger displacement for the arterial syringe at pass end. \\
\gui{End (venous)} & float & 0.0 & Plunger displacement for the venous syringe at pass end. \\
\bottomrule
\end{tabularx}

% --- 5.7 Advanced ---
\subsection{Advanced}

\begin{tabularx}{\textwidth}{L{3.5cm} l l X}
\toprule
\textbf{Parameter} & \textbf{Type} & \textbf{Default} & \textbf{Description} \\
\midrule
\gui{Custom G-code} & toggle & Off & Enable to provide custom G-code templates (reveals the Custom G-code tab). \\
\gui{Pathfinding algorithm} & select & DFS & DFS: modified depth-first search with KD-tree collision detection. \\
\gui{Overlap nodes} & int & 0 & Number of nodes to retrace at pass boundaries to bridge small gaps. \\
\gui{Overlap algorithm} & select & Retrace & ``Retrace'' scans all previous passes (original). ``Consecutive'' checks adjacent pairs (faster). \\
\gui{Generate plots} & toggle & On & Create interactive 3D HTML visualizations. \\
\gui{Speed calculation} & toggle & Off & Compute per-node print speed from vessel radius and flow rate. \\
\gui{Gap closure (SM)} & toggle & Off & Load gap extension files for SM prints. When on, shows file uploaders for pass indices and deltas. \\
\gui{Gap closure (MM)} & toggle & Off & Load gap extension files for MM prints. When on, shows file uploaders for pass indices and deltas. \\
\bottomrule
\end{tabularx}

% ===================================================================
\section{Pipeline Tab}
\label{sec:pipeline}
% ===================================================================

The \gui{Pipeline} tab is the primary workspace where you run the X-CAVATE
pipeline and view results.

\subsection{Running the Pipeline}

\begin{enumerate}[leftmargin=*]
    \item Upload both input files in the sidebar.
    \item Configure required and optional parameters.
    \item Click \btn{Run X-CAVATE} (full-width primary button at the top of the tab).
\end{enumerate}

\begin{infobox}
The run button is disabled until both input files are uploaded.
\end{infobox}

\subsection{Progress Bar}

The pipeline executes in 7 steps. A progress bar tracks each step:

\begin{tabularx}{\textwidth}{c X}
\toprule
\textbf{Step} & \textbf{Description} \\
\midrule
1 & Reading network and inlet/outlet files \\
2 & Interpolating network to nozzle resolution \\
3 & Building adjacency graph and detecting branchpoints \\
4 & Generating collision-free print passes \\
5 & Post-processing: subdivision, gap closure, overlap \\
6 & Processing multimaterial passes \\
7 & Writing output files and G-code \\
\bottomrule
\end{tabularx}

Upon completion, a success message appears: ``X-CAVATE finished successfully.''

\subsection{3D Visualizations}

If \gui{Generate plots} is enabled, three interactive 3D Plotly charts are
displayed below the progress bar:

\begin{enumerate}[leftmargin=*]
    \item \textbf{Original Network} --- The raw vascular network before processing, with each vessel shown as a separate trace.
    \item \textbf{Single Material Print Passes} --- The computed collision-free SM print paths, color-coded by pass index.
    \item \textbf{Multimaterial Print Passes} --- (If multimaterial is enabled) MM paths color-coded by material type (arterial vs.\ venous).
\end{enumerate}

All plots are interactive: rotate by clicking and dragging, zoom with the scroll
wheel, and hover over traces for coordinate information.  A download button
appears beneath each plot to save it as a standalone HTML file
(\code{original\_network.html}, \code{print\_passes\_SM.html},
\code{print\_passes\_MM.html}).

\subsection{Download Results}

Below the visualizations, three download columns appear:

\begin{tabularx}{\textwidth}{l X}
\toprule
\textbf{Column} & \textbf{Contents} \\
\midrule
\textbf{G-code} & Download buttons for each generated G-code file (e.g., \code{gcode\_SM\_pressure.txt}, \code{gcode\_MM\_pressure.txt}). \\
\textbf{Coordinate Files} & Download buttons for SM/MM coordinate files (x, y, z, combined). \\
\textbf{Changelog} & Download button for \code{changelog.txt} (gap closure operations log). \\
\bottomrule
\end{tabularx}

\begin{tipbox}
Results persist in the browser session. You can switch to other tabs and return to download files later without re-running the pipeline.
\end{tipbox}

% ===================================================================
\section{Custom G-code Tab}
\label{sec:custom-gcode}
% ===================================================================

This tab appears only when the \gui{Custom G-code} toggle is enabled in the
sidebar's Advanced section. It provides a tabbed editor for entering
printer-specific G-code snippets that replace the default commands in the
final output file.

\subsection{Sub-tabs}

The Custom G-code tab contains four sub-tabs:

\begin{tabularx}{\textwidth}{l X}
\toprule
\textbf{Sub-tab} & \textbf{What to paste} \\
\midrule
\gui{Header} & Printer initialization code: pressure box connections, homing commands, etc. \\
\gui{Single Material} & Start/stop extrusion code for single-nozzle prints (2 fields, side-by-side). \\
\gui{Multimaterial} & Start/stop extrusion + active/resting pressure for each printhead (8 fields in a 2-column layout). \\
\gui{Dwell} & Optional time-delay code for ``spot welding'' at pass junctions (2 fields: start-of-pass and end-of-pass). \\
\bottomrule
\end{tabularx}

\subsection{Template Fields}

The table below maps every editable text area to the output template file it
replaces:

\begin{longtable}{L{4.5cm} L{5.5cm} l}
\toprule
\textbf{GUI Label} & \textbf{Output Filename} & \textbf{Sub-tab} \\
\midrule
\endfirsthead
\toprule
\textbf{GUI Label} & \textbf{Output Filename} & \textbf{Sub-tab} \\
\midrule
\endhead
Header code & \code{header\_code.txt} & Header \\
Start extrusion & \code{start\_extrusion\_code.txt} & Single Material \\
Stop extrusion & \code{stop\_extrusion\_code.txt} & Single Material \\
Start extrusion --- Printhead 1 & \code{start\_extrusion\_code\_printhead1.txt} & Multimaterial \\
Start extrusion --- Printhead 2 & \code{start\_extrusion\_code\_printhead2.txt} & Multimaterial \\
Stop extrusion --- Printhead 1 & \code{stop\_extrusion\_code\_printhead1.txt} & Multimaterial \\
Stop extrusion --- Printhead 2 & \code{stop\_extrusion\_code\_printhead2.txt} & Multimaterial \\
Active pressure --- Printhead 1 & \code{active\_pressure\_printhead1.txt} & Multimaterial \\
Active pressure --- Printhead 2 & \code{active\_pressure\_printhead2.txt} & Multimaterial \\
Resting pressure --- Printhead 1 & \code{rest\_pressure\_printhead1.txt} & Multimaterial \\
Resting pressure --- Printhead 2 & \code{rest\_pressure\_printhead2.txt} & Multimaterial \\
Dwell code (start of pass) & \code{dwell\_start.txt} & Dwell \\
Dwell code (end of pass) & \code{dwell\_end.txt} & Dwell \\
\bottomrule
\end{longtable}

\begin{tipbox}
Snippets are saved automatically in the browser session and persist across
page re-runs, even when the Custom G-code toggle is temporarily turned off.
\end{tipbox}

% ===================================================================
\section{Print Instructions Tab}
\label{sec:instructions}
% ===================================================================

After a successful pipeline run, the \gui{Print Instructions} tab shows
positioning and printing guidance specific to your network and parameters.

\subsection{Start Position Metrics}

Three metric cards display the starting coordinates:

\begin{tabularx}{\textwidth}{l X}
\toprule
\textbf{Metric} & \textbf{Description} \\
\midrule
\gui{Start X} (mm) & X-coordinate where the nozzle should begin. \\
\gui{Start Y} (mm) & Y-coordinate where the nozzle should begin. \\
\gui{Start Z} (mm) & Z-coordinate where the nozzle should begin. \\
\bottomrule
\end{tabularx}

\subsection{Padding Metrics}

Four metric cards show the container padding --- how far the network is offset
from each container wall:

\begin{tabularx}{\textwidth}{l X}
\toprule
\textbf{Metric} & \textbf{Description} \\
\midrule
\gui{Left padding} & Distance (mm) from the network to the left container wall. \\
\gui{Right padding} & Distance (mm) from the network to the right container wall. \\
\gui{Front padding} & Distance (mm) from the network to the front container wall. \\
\gui{Back padding} & Distance (mm) from the network to the back container wall. \\
\bottomrule
\end{tabularx}

\subsection{Printing Instructions}

Depending on the mode, one of the following expandable sections appears:

\begin{description}[style=nextline]
    \item[\gui{Single Material Instructions}]
    Step-by-step instructions for positioning the single printhead and starting
    the print job.
    \item[\gui{Multimaterial Calibration}]
    Instructions for calibrating the offset between the two printheads.
    \item[\gui{Multimaterial Positioning}]
    Instructions for positioning both printheads and starting the multimaterial
    print job.
\end{description}

\begin{infobox}
This tab shows ``Run the pipeline first to see print instructions'' until you
have completed a successful pipeline run.
\end{infobox}

% ===================================================================
\section{Calibration Tab}
\label{sec:calibration}
% ===================================================================

The \gui{Calibration} tab provides two features: computing the volumetric flow
rate~($Q$) from calibration measurements and validating measured filament
diameters against target values.

\subsection{Compute Volumetric Flow Rate (Q)}

This section calculates the average volumetric flow rate from calibration
filament measurements and produces a calibration curve (radius vs.\ speed).
During calibration, pressure is held constant while the print speed is varied
for each target diameter group.

\subsubsection{Inputs}

\begin{tabularx}{\textwidth}{L{3.5cm} l X}
\toprule
\textbf{Widget} & \textbf{Type} & \textbf{Description} \\
\midrule
\gui{Upload measured diameters CSV} & file upload & CSV file with a \code{Length} column containing measured filament diameters in \textbf{microns}. \\
\gui{Print speeds used during calibration} & text input & Comma-separated list of print speeds in mm/s, one per group. Pressure is held constant; speed is varied for each calibration filament. \\
\gui{Measurements per speed} & int (default: 3) & Number of rows in the CSV to average for each print speed group. \\
\bottomrule
\end{tabularx}

\subsubsection{Running Q Computation}

Click \btn{Compute Q} to process. The button is disabled until a CSV file is
uploaded.

\subsubsection{Outputs}

Upon completion, the section displays:

\begin{enumerate}[leftmargin=*]
    \item The computed $Q$ value as a metric (mm\textsuperscript{3}/s).
    \item An interactive Plotly chart showing the calibration curve (radius vs.\ speed).
    \item Download buttons for:
    \begin{itemize}[nosep]
        \item \textbf{Q\_pressure.txt} --- the computed $Q$ value.
        \item \textbf{radii\_pressure.html} --- interactive calibration curve.
    \end{itemize}
\end{enumerate}

\subsection{Calibration Validation}

This section lets you compare measured filament diameters against target values
to verify print accuracy.

\subsubsection{Inputs}

\begin{tabularx}{\textwidth}{L{3.5cm} l X}
\toprule
\textbf{Widget} & \textbf{Type} & \textbf{Description} \\
\midrule
\gui{Upload measured diameters CSV} & file upload & CSV file with a \code{Length} column containing measured filament diameters in \textbf{microns}. \\
\gui{Target diameters} & text input & Comma-separated list of expected filament diameters in \textbf{mm}, in the same order as the measurements. \\
\gui{Measurements per target} & int (default: 3) & Number of rows in the CSV to average for each target diameter. \\
\bottomrule
\end{tabularx}

\subsubsection{Running Validation}

Click \btn{Run Validation} to process. The button is disabled until a CSV file
is uploaded.

\subsubsection{Outputs}

Upon completion, the section displays:

\begin{enumerate}[leftmargin=*]
    \item An interactive Plotly chart comparing target vs.\ measured diameters.
    \item Download buttons for:
    \begin{itemize}[nosep]
        \item \textbf{CSV report} (\code{calibration\_error.csv}) --- tabular error analysis.
        \item \textbf{HTML report} (\code{calibration\_error.html}) --- interactive visualization.
    \end{itemize}
\end{enumerate}

% ===================================================================
\section{Output Files}
\label{sec:outputs}
% ===================================================================

All output is written to a temporary directory. Use the download buttons in the
\gui{Pipeline} tab to save files locally. When running via CLI, files are
written to the \code{outputs/} directory (or the path specified by
\code{-{}-output\_dir}).

\subsection{Directory Structure}

\begin{lstlisting}
outputs/
  graph/
    graph.txt               # Adjacency graph
    special_nodes.txt       # Branchpoints, endpoints, inlets, outlets
    SM_x_coords.txt         # Single-material x coordinates per pass
    SM_y_coords.txt         # Single-material y coordinates per pass
    SM_z_coords.txt         # Single-material z coordinates per pass
    SM_combined.txt         # Combined xyz per pass
    MM_x_coords.txt         # Multimaterial x coordinates (if enabled)
    MM_y_coords.txt         # Multimaterial y coordinates (if enabled)
    MM_z_coords.txt         # Multimaterial z coordinates (if enabled)
    MM_combined.txt         # Multimaterial combined xyz (if enabled)
  gcode/
    gcode_SM_pressure.txt   # Single-material G-code
    gcode_MM_pressure.txt   # Multimaterial G-code (if enabled)
  plots/
    network_original.html   # 3D plot of original network
    network_SM.html         # 3D plot of SM print passes
    network_MM.html         # 3D plot of MM print passes (if enabled)
  changelog.txt             # Gap closure operations log
\end{lstlisting}

\begin{infobox}
The G-code filename suffix matches your selected printer type:
\code{pressure}, \code{positive\_ink}, or \code{aerotech}.
\end{infobox}

\subsection{File Descriptions}

\begin{longtable}{L{4cm} L{2cm} X}
\toprule
\textbf{File} & \textbf{Directory} & \textbf{Description} \\
\midrule
\endfirsthead
\toprule
\textbf{File} & \textbf{Directory} & \textbf{Description} \\
\midrule
\endhead
\code{graph.txt} & \code{graph/} & Full adjacency graph of the interpolated network. \\
\code{special\_nodes.txt} & \code{graph/} & Identified branchpoints, endpoints, inlets, and outlets. \\
\code{SM\_x\_coords.txt} & \code{graph/} & X-coordinates for each SM print pass, one pass per line. \\
\code{SM\_y\_coords.txt} & \code{graph/} & Y-coordinates for each SM print pass. \\
\code{SM\_z\_coords.txt} & \code{graph/} & Z-coordinates for each SM print pass. \\
\code{SM\_combined.txt} & \code{graph/} & Combined (x, y, z) coordinates per SM pass. \\
\code{MM\_*.txt} & \code{graph/} & Multimaterial equivalents of the above (if enabled). \\
\code{gcode\_SM\_*.txt} & \code{gcode/} & Single-material G-code for the selected printer type. \\
\code{gcode\_MM\_*.txt} & \code{gcode/} & Multimaterial G-code (if enabled). \\
\code{network\_original.html} & \code{plots/} & Interactive 3D visualization of the original vascular network. \\
\code{network\_SM.html} & \code{plots/} & Interactive 3D visualization of SM print passes. \\
\code{network\_MM.html} & \code{plots/} & Interactive 3D visualization of MM print passes. \\
\code{changelog.txt} & root & Log of all gap closure operations performed. \\
\bottomrule
\end{longtable}

% ===================================================================
\section{Troubleshooting}
\label{sec:troubleshooting}
% ===================================================================

\begin{longtable}{L{4.5cm} X}
\toprule
\textbf{Problem} & \textbf{Solution} \\
\midrule
\endfirsthead
\toprule
\textbf{Problem} & \textbf{Solution} \\
\midrule
\endhead
\code{Runtime instance already exists} (GUI) &
A stale Streamlit process is running. Kill it with \code{pkill -f streamlit}, then restart the server. \\
\midrule
Inlet/outlet matching warnings &
Inlet/outlet coordinates are approximate. X-CAVATE matches to the nearest network point and logs the distance. Verify the matches are reasonable. \\
\midrule
Pipeline hangs on large networks &
Reduce coordinates per vessel to 100 or fewer, or enable downsampling. \\
\midrule
\code{ModuleNotFoundError} &
Run \code{pip install ".[gui]"} to install all dependencies, or recreate the conda environment. \\
\midrule
Docker container won't start &
Ensure Docker Desktop is running and port 8501 is not in use by another process. \\
\midrule
Download buttons disappear after clicking &
Update to the latest version of X-CAVATE (this has been fixed). \\
\midrule
No G-code files generated &
Verify that both input files are correctly formatted and that the pipeline completed without errors. Check the terminal for error messages. \\
\midrule
Plots not showing &
Ensure \gui{Generate plots} is enabled in the Advanced section. If disabled, plots are skipped for faster execution. \\
\midrule
Browser does not open &
Navigate manually to \code{http://localhost:8501}. On remote servers, ensure port 8501 is forwarded. \\
\bottomrule
\end{longtable}

\end{document}
